\chapter{Experiments and Results}
\label{chap:experiments_results}

\section{Overview}
\label{sec:results_overview}

This chapter evaluates the Block-Insure prototype against the research questions
defined in Chapter~\ref{chap:methodology}. The evaluation covers smart-contract
correctness, code coverage, deployability constraints, gas efficiency of Merkle
commitment, synthetic fraud-risk classification, and exploratory local
claim-submission throughput. The auditor-reputation study is also reported as an
implemented but currently unmeasurable study because no finalized voting
sessions exist.

The results must be interpreted within the experimental boundary. The
blockchain is a local Hardhat network, healthcare records are deterministic and
synthetic, both oracle workers share prototype infrastructure, and the
throughput study consists of one local batch at each concurrency level.
Consequently, the results demonstrate implemented behavior and reproducibility;
they do not establish production performance, public-chain deployability, or
real-world fraud-model generalization.

\begin{table}[htbp]
\centering
\caption{Summary of research questions and principal findings}
\label{tab:rq_findings}
\resizebox{\textwidth}{!}{
\begin{tabular}{|p{1.1cm}|p{7.1cm}|p{7.0cm}|}
\hline
\textbf{RQ} & \textbf{Principal Finding} & \textbf{Interpretation} \\
\hline
RQ1 & 117 contract tests passed; statement, function, and line coverage exceeded
98\%, while branch coverage was 77.21\% & Strong evidence of conformance to the
specified test suite, but not a security audit \\
\hline
RQ2 & Merkle-root storage saved 82.41\%--99.62\% gas across tested registry
sizes & Compact commitment scales substantially better than individual hash
storage \\
\hline
RQ3 & Held-out synthetic evaluation produced \(TP=7\), \(TN=8\), \(FP=0\),
\(FN=0\), with AUC and AP of 1.0000 & The model separates deliberately
structured synthetic cases; result is not generalizable to real fraud \\
\hline
RQ4 & Throughput increased to 48.3827 successful claims/s at concurrency 20
with no failures; concurrency 50 achieved 49.2880 successful claims/s but only
78\% success & Local automining and nonce handling became the limiting factor \\
\hline
RQ5 & No finalized voting sessions were available & No correlation is reported;
fabricating a result was avoided \\
\hline
\end{tabular}
}
\end{table}

\section{Experimental Environment and Verification Status}
\label{sec:experimental_environment}

The implementation and generated artifacts were audited on 4 June 2026 using
the local workspace configuration shown in Table~\ref{tab:environment}.
Contract tests, backend analytics tests, oracle syntax validation, frontend lint,
and the frontend build attempt were rerun during the audit. The retained
coverage report was generated earlier by the same project test suite.
Node-visible memory is reported because it was measured programmatically;
installed physical memory should be manually confirmed before final thesis
submission.

\begin{table}[htbp]
\centering
\caption{Experimental environment}
\label{tab:environment}
\begin{tabular}{|p{5.0cm}|p{8.4cm}|}
\hline
\textbf{Component} & \textbf{Observed Configuration} \\
\hline
Operating system & Windows, \texttt{win32 x64}, release \texttt{10.0.26200} \\
\hline
Processor & AMD Ryzen 5 5600H with Radeon Graphics, 12 logical processors \\
\hline
Node-visible memory & 7,883,206,656 bytes, approximately 7.34 GiB \\
\hline
Node.js and npm & Node.js \texttt{v24.16.0}; npm \texttt{10.2.5} \\
\hline
Blockchain toolchain & Solidity \texttt{0.8.28}; Hardhat \texttt{2.28.0};
OpenZeppelin Contracts \texttt{5.6.1}; ethers.js \texttt{6.16.0} \\
\hline
Application stack & Express \texttt{5.2.1}; Mongoose \texttt{9.6.2}; React
\texttt{19.2.6}; Vite \texttt{8.0.14} \\
\hline
Blockchain network & Local Hardhat network with automining and unlimited
contract size enabled \\
\hline
\end{tabular}
\end{table}

\begin{table}[htbp]
\centering
\caption{Verification and artifact status}
\label{tab:verification_status}
\begin{tabular}{|p{6.0cm}|p{3.2cm}|p{4.4cm}|}
\hline
\textbf{Check or Study} & \textbf{Status} & \textbf{Observed Result} \\
\hline
Hardhat smart-contract tests & Complete & 117 passing, 0 failing \\
\hline
Solidity coverage artifact & Complete & 98.64\% statements; 77.21\% branches;
100\% functions; 98.91\% lines \\
\hline
Backend analytics tests & Complete & Passed \\
\hline
Oracle syntax validation & Complete & Passed \\
\hline
Frontend lint & Complete & Passed \\
\hline
Frontend production build & Environment-blocked & Vite/Rolldown rejected the
Windows workspace path containing spaces \\
\hline
Gas comparison & Complete & Six registry sizes measured \\
\hline
Fraud-risk evaluation and chart generation & Complete & Held-out artifacts and
10 PNG charts generated \\
\hline
Claim throughput benchmark & Complete & Five concurrency levels measured \\
\hline
Auditor reputation analysis & Not yet measurable & No finalized voting
sessions \\
\hline
\end{tabular}
\end{table}

\section{Smart-Contract Correctness and Coverage}
\label{sec:contract_results}

\subsection{Functional Test Results}

The Hardhat suite exercised successful operations, invalid inputs, unauthorized
calls, duplicate detection, oracle quorum, timeout behavior, appeals, voting,
settlement, closure, and emergency controls. All 117 tests passed.

\begin{table}[htbp]
\centering
\caption{Smart-contract test-suite results}
\label{tab:contract_tests}
\begin{tabular}{|p{8.0cm}|r|p{2.8cm}|}
\hline
\textbf{Test Category} & \textbf{Passing Tests} & \textbf{Result} \\
\hline
Admin decision and settlement & 20 & Passed \\
\hline
Claim appeal workflow & 5 & Passed \\
\hline
Claim submission & 13 & Passed \\
\hline
Fraud and duplicate detection & 7 & Passed \\
\hline
Roles, pause, and voting & 22 & Passed \\
\hline
Oracle contract logic & 16 & Passed \\
\hline
Policy package management & 7 & Passed \\
\hline
Policy purchase & 12 & Passed \\
\hline
Registry Merkle-root commitment & 2 & Passed \\
\hline
On-chain risk score & 5 & Passed \\
\hline
System-design hardening & 8 & Passed \\
\hline
\textbf{Total} & \textbf{117} & \textbf{Passed} \\
\hline
\end{tabular}
\end{table}

The result supports RQ1 by showing that the implemented contract conforms to the
tested lifecycle and authorization requirements. Passing tests do not establish
the absence of vulnerabilities outside the test cases. No independent audit,
formal verification, invariant-based fuzzing, or public-network test is claimed.
The risk-score tests also confirm that a clean claim receives 90 before oracle
verification. Since the score's conditions overlap submission validation and
duplicate checks, this should be interpreted as lifecycle confidence rather
than empirical fraud-risk discrimination.

\subsection{Coverage Results}

\begin{table}[htbp]
\centering
\caption{Solidity coverage results}
\label{tab:coverage_results}
\begin{tabular}{|p{4.0cm}|r|r|r|}
\hline
\textbf{Coverage Type} & \textbf{Covered} & \textbf{Total} &
\textbf{Coverage (\%)} \\
\hline
Statements & 217 & 220 & 98.64 \\
\hline
Branches & 210 & 272 & 77.21 \\
\hline
Functions & 50 & 50 & 100.00 \\
\hline
Lines & 273 & 276 & 98.91 \\
\hline
\end{tabular}
\end{table}

Statement, function, and line coverage are high. Branch coverage is materially
lower at 77.21\%, indicating that some combinations of conditions and revert
paths remain unexecuted. The branch result is therefore a more informative
remaining testing gap than the overall passing-test count.

\subsection{Deployability Constraint}

\begin{table}[htbp]
\centering
\caption{Compiled contract size compared with EIP-170}
\label{tab:contract_size}
\begin{tabular}{|p{6.0cm}|r|}
\hline
\textbf{Measurement} & \textbf{Bytes} \\
\hline
Creation bytecode & 24,888 \\
\hline
Deployed runtime bytecode & 24,340 \\
\hline
EIP-170 deployed-bytecode limit & 24,576 \\
\hline
Remaining margin below limit & 236 \\
\hline
Percentage below limit & 0.96\% \\
\hline
\end{tabular}
\end{table}

The optimized contract is deployable on an EIP-170-enforcing EVM network.
Hardhat's unlimited-size bypass is disabled, and a regression test enforces the
24,576-byte boundary. The small remaining margin means future additions must
replace existing code or use governed modules rather than growing the monolith.

\section{Gas-Efficiency Study}
\label{sec:gas_results}

\subsection{Protocol}

The gas study compared two first-write operations on a fresh
\texttt{GasTestContract}: storing \(N\) individual \texttt{bytes32} hashes in a
mapping and storing one \texttt{bytes32} Merkle root. Receipt gas includes
transaction execution and calldata for the tested operation. The independent
variable was
\begin{equation}
N\in\{10,25,50,100,200,500\}.
\end{equation}
For each \(N\),
\begin{equation}
Saving(\%)=
\frac{Gas_{\mathrm{individual}}-Gas_{\mathrm{root}}}
{Gas_{\mathrm{individual}}}\times100.
\label{eq:gas_saving_result}
\end{equation}

\subsection{Measured Results}

\begin{table}[htbp]
\centering
\caption{Gas comparison: individual hashes versus one Merkle root}
\label{tab:gas_results}
\begin{tabular}{|r|r|r|r|r|}
\hline
\textbf{Records} & \textbf{Individual Gas} & \textbf{Root Gas} &
\textbf{Gas Saved} & \textbf{Saving (\%)} \\
\hline
10 & 249,300 & 43,852 & 205,448 & 82.41 \\
\hline
25 & 590,754 & 43,852 & 546,902 & 92.58 \\
\hline
50 & 1,159,916 & 43,852 & 1,116,064 & 96.22 \\
\hline
100 & 2,298,072 & 43,852 & 2,254,220 & 98.09 \\
\hline
200 & 4,574,588 & 43,852 & 4,530,736 & 99.04 \\
\hline
500 & 11,403,692 & 43,852 & 11,359,840 & 99.62 \\
\hline
\end{tabular}
\end{table}

Across these six observations, a least-squares linear fit for individual storage
is approximately
\begin{equation}
Gas_{\mathrm{individual}}(N)\approx21,686+22,764N,
\end{equation}
whereas first-write root storage remained 43,852 gas. The result supports RQ2:
compact commitment is substantially more scalable for anchoring a registry
snapshot. It does not include off-chain tree-construction cost, proof-generation
cost, or the gas cost of verifying inclusion proofs on-chain because the current
prototype verifies proofs off-chain.

\section{Bayesian Fraud-Risk Evaluation}
\label{sec:risk_results}

\subsection{Dataset and Evaluation Protocol}

The deterministic synthetic registry contains 72 records: 39 fraud-labeled and
33 legitimate. After stable sorting by hospital and invoice, every fifth record
was held out and the remaining four were used for training. This produced 57
training records and 15 held-out records.

\begin{table}[htbp]
\centering
\caption{Synthetic dataset composition}
\label{tab:dataset_composition}
\begin{tabular}{|p{6.5cm}|r|r|r|}
\hline
\textbf{Split} & \textbf{Total} & \textbf{Fraud} & \textbf{Legitimate} \\
\hline
Full registry & 72 & 39 & 33 \\
\hline
Training & 57 & 32 & 25 \\
\hline
Held-out & 15 & 7 & 8 \\
\hline
\end{tabular}
\end{table}

Model likelihoods and anomaly statistics were calculated only from the training
split. Thresholds from 0 to 100 in steps of five were evaluated on the training
records. The selected threshold maximized training F1, with ties broken by
higher precision and then the higher threshold. The selected threshold,
\(\theta=85\), was applied once to the held-out records. Held-out labels were
not used for threshold selection.

\subsection{Threshold Selection}

\begin{table}[htbp]
\centering
\caption{Representative training threshold-selection results}
\label{tab:threshold_selection_results}
\begin{tabular}{|r|r|r|r|r|}
\hline
\textbf{Threshold} & \textbf{Precision} & \textbf{Recall} &
\textbf{F1} & \textbf{Selection Role} \\
\hline
100 & 1.0000 & 0.5313 & 0.6939 & Too restrictive \\
\hline
95 & 1.0000 & 0.8438 & 0.9153 & High precision, lower recall \\
\hline
90 & 1.0000 & 0.8438 & 0.9153 & High precision, lower recall \\
\hline
85 & 1.0000 & 1.0000 & 1.0000 & Selected threshold \\
\hline
20 & 1.0000 & 1.0000 & 1.0000 & Tied maximum F1 \\
\hline
5 & 1.0000 & 1.0000 & 1.0000 & Tied maximum F1 \\
\hline
0 & 0.5614 & 1.0000 & 0.7191 & Predicts every training record as fraud \\
\hline
\end{tabular}
\end{table}

The broad interval of perfect training F1, from threshold 5 through 85,
demonstrates unusually strong separation in the synthetic data. It also shows
that the selected value of 85 is primarily a consequence of the explicit
tie-breaking rule, not evidence that 85 is a generally optimal real-world fraud
threshold.

\subsection{Held-Out Classification}

\begin{table}[htbp]
\centering
\caption{Held-out confusion matrix at \(\theta=85\)}
\label{tab:heldout_confusion}
\begin{tabular}{|p{4.3cm}|c|c|}
\hline
\textbf{Actual / Predicted} & \textbf{Fraud} & \textbf{Legitimate} \\
\hline
Actual Fraud & \(TP=7\) & \(FN=0\) \\
\hline
Actual Legitimate & \(FP=0\) & \(TN=8\) \\
\hline
\end{tabular}
\end{table}

\begin{table}[htbp]
\centering
\caption{Held-out classification metrics}
\label{tab:heldout_metrics}
\begin{tabular}{|p{6.5cm}|r|}
\hline
\textbf{Metric} & \textbf{Value} \\
\hline
Accuracy & 1.0000 \\
\hline
Precision & 1.0000 \\
\hline
Recall / sensitivity & 1.0000 \\
\hline
Specificity & 1.0000 \\
\hline
False-positive rate & 0.0000 \\
\hline
F1 score & 1.0000 \\
\hline
ROC AUC & 1.0000 \\
\hline
Average Precision & 1.0000 \\
\hline
Fraud prevalence & 0.4667 \\
\hline
Average held-out risk score & 49.4 \\
\hline
\end{tabular}
\end{table}

\subsection{Baseline Comparison}

\begin{table}[htbp]
\centering
\caption{Held-out comparison with simple baselines}
\label{tab:baseline_results}
\resizebox{\textwidth}{!}{
\begin{tabular}{|p{4.7cm}|r|r|r|r|r|r|r|r|}
\hline
\textbf{Model} & \textbf{TP} & \textbf{TN} & \textbf{FP} & \textbf{FN} &
\textbf{Accuracy} & \textbf{Precision} & \textbf{Recall} & \textbf{F1} \\
\hline
Bayesian model, \(\theta=85\) & 7 & 8 & 0 & 0 & 1.0000 & 1.0000 &
1.0000 & 1.0000 \\
\hline
Always Fraud & 7 & 0 & 8 & 0 & 0.4667 & 0.4667 & 1.0000 & 0.6364 \\
\hline
Amount above training mean (0.330175 ETH) & 3 & 6 & 2 & 4 & 0.6000 &
0.6000 & 0.4286 & 0.5000 \\
\hline
\end{tabular}
}
\end{table}

The Bayesian model outperformed both simple baselines on the held-out synthetic
set. The comparison demonstrates that the result is not explained merely by
class prevalence or claim amount.

\subsection{Fraud-Scenario Breakdown}

\begin{table}[htbp]
\centering
\caption{Held-out risk score by synthetic scenario}
\label{tab:fraud_scenario_breakdown}
\begin{tabular}{|p{5.0cm}|r|r|r|}
\hline
\textbf{Scenario} & \textbf{Records} & \textbf{Detection Rate} &
\textbf{Mean Risk Score} \\
\hline
Legitimate & 8 & 0.0000 & 5.5 \\
\hline
Used invoice & 1 & 1.0000 & 100.0 \\
\hline
Date mismatch & 1 & 1.0000 & 100.0 \\
\hline
Suspicious repeated pattern & 2 & 1.0000 & 100.0 \\
\hline
Cancelled record & 1 & 1.0000 & 99.0 \\
\hline
Blacklisted hospital & 2 & 1.0000 & 99.0 \\
\hline
\end{tabular}
\end{table}

The scenario table explains the perfect aggregate metrics. Fraud cases were
constructed with strong indicators such as invalid registry status, repeated
claim count, cancelled invoice, suspended license, or blacklisted hospital,
while most legitimate records activated only the clean-match factor. One
legitimate record received a manual-review recommendation due to a bill-range
anomaly but remained below the fraud threshold.

\begin{figure}[htbp]
\centering
\includegraphics[width=0.75\textwidth]
{backend/evaluation-results/evaluation-charts/confusion_matrix.png}
\caption{Held-out confusion matrix for the Bayesian fraud model}
\label{fig:confusion_matrix}
\end{figure}

\begin{figure}[htbp]
\centering
\includegraphics[width=0.75\textwidth]
{backend/evaluation-results/evaluation-charts/roc_curve.png}
\caption{Held-out ROC curve}
\label{fig:roc_curve}
\end{figure}

\begin{figure}[htbp]
\centering
\includegraphics[width=0.75\textwidth]
{backend/evaluation-results/evaluation-charts/precision_recall_curve.png}
\caption{Held-out precision-recall curve}
\label{fig:precision_recall_curve}
\end{figure}

\subsection{Interpretation}

The results support RQ3 only within the synthetic experiment. The training and
held-out protocol avoids direct test-label leakage during parameter training and
threshold selection. However, the data-generation rules strongly couple fraud
labels to observable statuses and signals. Evaluation queries are also
constructed from the same registry records, so ordinary legitimate cases are
exact matches. Furthermore, the runtime rule-based oracle comparison directly
checks the synthetic \texttt{fraudLabel}, although that label is excluded from
the Bayesian evidence factors. The perfect AUC, AP, and confusion matrix
therefore demonstrate correct separation of designed scenarios, not expected
performance on noisy or adversarial real insurance data.

\section{Claim-Submission Throughput Study}
\label{sec:throughput_results}

\subsection{Protocol and Measurement Boundary}

The exploratory benchmark tested
\begin{equation}
C\in\{1,5,10,20,50\}
\end{equation}
concurrent claim submissions. Before each timed batch, the script purchased one
policy per claim; setup time was excluded. Every claim used unique document and
invoice hashes. The script measured:
\begin{enumerate}
    \item response time of the backend \texttt{/health} endpoint;
    \item blockchain \texttt{submitClaim} send-to-receipt time;
    \item combined end-to-end time; and
    \item successful-claim throughput based on batch wall-clock duration.
\end{enumerate}

The backend measurement represents basic API responsiveness, not the latency of
a backend claim-processing endpoint. All blockchain submissions used one signer
managed by an ethers.js \texttt{NonceManager}. Each concurrency level was run
once; therefore, the results are exploratory rather than statistically stable.

\subsection{Results}

\begin{table}[htbp]
\centering
\caption{Local claim-submission throughput and latency}
\label{tab:throughput_results}
\resizebox{\textwidth}{!}{
\begin{tabular}{|r|r|r|r|r|r|r|r|r|}
\hline
\textbf{Concurrency} & \textbf{Successful} & \textbf{Failed} &
\textbf{Success (\%)} & \textbf{Avg Backend (ms)} &
\textbf{Avg Blockchain (ms)} & \textbf{Avg E2E (ms)} &
\textbf{P95 E2E (ms)} & \textbf{Successful Claims/s} \\
\hline
1 & 1 & 0 & 100 & 2.78 & 125.28 & 151.99 & 151.99 & 6.5740 \\
\hline
5 & 5 & 0 & 100 & 3.88 & 201.61 & 219.06 & 219.76 & 22.7484 \\
\hline
10 & 10 & 0 & 100 & 11.59 & 293.08 & 332.19 & 334.39 & 29.8951 \\
\hline
20 & 20 & 0 & 100 & 14.65 & 374.95 & 411.42 & 412.74 & 48.3827 \\
\hline
50 & 39 & 11 & 78 & 46.54 & 678.55 & 787.85 & 790.66 & 49.2880 \\
\hline
\end{tabular}
}
\end{table}

Throughput increased rapidly up to concurrency 20, where all submissions
succeeded at 48.3827 claims/s. Increasing concurrency to 50 produced only a
small throughput gain, to 49.2880 successful claims/s, while success rate fell
to 78\%. All 11 failures reported ``nonce too high'' because local Hardhat
automining does not queue out-of-order transactions. The result therefore
identifies a local signer/automining saturation point, not a smart-contract
validation limit.

Blockchain confirmation dominated the measured end-to-end duration. Depending
on concurrency, average blockchain time represented approximately
82.43\%--92.03\% of average end-to-end time, while the health-endpoint response
represented approximately 1.77\%--5.91\%. This supports RQ4's local conclusion
that transaction processing, rather than basic backend responsiveness, is the
primary measured bottleneck.

\begin{figure}[htbp]
\centering
\includegraphics[width=0.82\textwidth]
{backend/evaluation-results/evaluation-charts/throughput_vs_latency.png}
\caption{Local successful throughput and latency by concurrency}
\label{fig:throughput_latency}
\end{figure}

\section{Auditor-Reputation Study Status}
\label{sec:auditor_result_status}

The implemented analysis groups finalized votes by auditor, calculates
historical agreement with finalized consensus, reads current on-chain
reputation, and computes Pearson correlation:
\begin{equation}
r=
\frac{\sum_i(x_i-\bar{x})(y_i-\bar{y})}
{\sqrt{\sum_i(x_i-\bar{x})^2\sum_i(y_i-\bar{y})^2}}.
\label{eq:auditor_pearson_result}
\end{equation}
At least two auditors and nonzero variation are required. No finalized voting
sessions existed at the time of evaluation; consequently, no defensible
correlation can be calculated. RQ5 therefore remains unanswered. Agreement with
consensus would in any case be only a proxy for correctness because consensus
itself is not independent ground truth.

\section{Threats to Validity}
\label{sec:threats_to_validity}

\begin{table}[htbp]
\centering
\caption{Principal threats to validity}
\label{tab:validity_threats}
\resizebox{\textwidth}{!}{
\begin{tabular}{|p{3.3cm}|p{6.3cm}|p{6.0cm}|}
\hline
\textbf{Validity Type} & \textbf{Threat} & \textbf{Mitigation or Required Future
Work} \\
\hline
Internal validity & Deterministic synthetic generation creates strongly
separable fraud patterns; throughput has one batch per level & Repeat with
multiple random seeds, repeated trials, confidence intervals, and independent
datasets \\
\hline
Construct validity & Passing tests measure conformance to specified behavior,
not security; health-endpoint latency is not claim-API latency & Add independent
security audit, fuzzing, full workflow latency, and public-network experiments \\
\hline
External validity & Local Hardhat automining, one machine, synthetic healthcare
records, and shared oracle infrastructure differ from production & Deploy a
modular contract to a public testnet and use independently operated data sources \\
\hline
Conclusion validity & Held-out set contains only 15 records; AUC and AP are
perfect because of strong separation; no repeated throughput runs & Use larger
expert-labeled data, cross-validation, stronger baselines, and statistical
uncertainty estimates \\
\hline
Security validity & Merkle-root matching, evidence-chain integrity, and reserve
intelligence are audit signals rather than fully enforced gates & Anchor
evidence-chain heads, require committed-root proof validity, constrain
withdrawals, and use multisignature administration \\
\hline
\end{tabular}
}
\end{table}

\section{Result Synthesis}
\label{sec:result_synthesis}

\begin{table}[htbp]
\centering
\caption{Implemented objective mapped to experimental evidence}
\label{tab:objective_evidence}
\resizebox{\textwidth}{!}{
\begin{tabular}{|p{5.7cm}|p{7.1cm}|p{3.0cm}|}
\hline
\textbf{Objective} & \textbf{Evidence} & \textbf{Assessment} \\
\hline
Enforce auditable lifecycle and role boundaries & 117 passing contract tests;
98.64\% statement and 98.91\% line coverage & Demonstrated within tested scope \\
\hline
Detect exact duplicate claims & Passing document-hash, invoice-hash, and
wallet/date/type duplicate tests & Demonstrated for exact duplicates \\
\hline
Provide fault-aware external verification & Two-confirmation quorum, split
decision fails conservatively, timeout tested & Demonstrated in local simulation \\
\hline
Reduce registry commitment gas & 82.41\%--99.62\% gas savings & Demonstrated for
tested first-write operations \\
\hline
Evaluate explainable fraud scoring & Held-out metrics, baselines, AUC, AP, and
scenario breakdown & Demonstrated only on synthetic cases \\
\hline
Measure concurrent submissions & Five local concurrency levels; saturation and
nonce failures identified & Exploratory local evidence \\
\hline
Validate auditor reputation & Analysis tooling implemented & Pending data \\
\hline
Support public-network deployment & Runtime bytecode is 236 bytes below EIP-170
& Achieved for Cancun-compatible EVM networks \\
\hline
\end{tabular}
}
\end{table}

\section{Summary}

The evaluation demonstrates that Block-Insure implements and tests a substantial
claim-processing lifecycle, produces compact and highly gas-efficient registry
commitments, and separates deliberately designed synthetic fraud scenarios. It
also exposes important limitations that improve the credibility of the result:
branch coverage remains lower than other coverage measures, contract-size
headroom remains narrow, fraud metrics are based on a small deterministic dataset,
high-concurrency results are affected by local nonce handling, and the auditor
study has no data. The defensible conclusion is that Block-Insure is a working
and reproducible research prototype with clearly identified paths toward
production-grade decentralization, privacy, validation, and deployment.
